% HARNAIS-HORS-SECTION: guide de lecture contenant des renvois de chapitres et de sections,
% sans nouvelle valeur empirique ni nouveau parametre.
\begin{table}[H]\centering\small
\begingroup
\definecolor{reperecritique}{HTML}{A52D39}
\definecolor{reperereponse}{HTML}{14828B}
\renewcommand{\arraystretch}{1.3}
\begin{tabularx}{\textwidth}{@{}*{3}{>{\raggedright\arraybackslash}X}@{}}
\toprule
\textbf{Question à résoudre} & \textcolor{reperecritique}{\textbf{Méthode et apport}} & \textcolor{reperereponse}{\textbf{Résultat à examiner et limite}} \\
\midrule
Comment représenter l'effet d'un état de conformité que le forfait ne distingue pas ?
& Relier un état supposé aux paramètres d'une même chaîne de pertes (chapitres~\ref{chap:socle} à~\ref{chap:multietats}).
& Comparer les capitaux annuels (\S\ref{sec:chiffre-de-tete}) ; les écarts entre états sont posés, sans effet causal établi. \\
\addlinespace
Que reste-t-il du calcul si les données ne donnent pas le sens des transmissions ?
& Examiner les sources, séparer force et asymétrie, puis explorer les directions admissibles (chapitres~\ref{ch:identifiabilite} et~\ref{chap:partielle}).
& Une plage aux sommets, à force commune fixée ; ni intervalle de confiance ni borne globale démontrée. \\
\addlinespace
Quelle correction comparer, et quelles données manquent pour décider ?
& Calculer les effets de fermeture, puis éprouver les coûts et la calibration (chapitres~\ref{chap:resultats} et~\ref{ch:robustesse}).
& Gains conditionnels et besoins de collecte ; l'efficacité réelle d'une remédiation reste à mesurer. \\
\bottomrule
\end{tabularx}
\endgroup
\caption{Les questions qui guident le mémoire et les réponses à vérifier. Les renvois indiquent
où la méthode est construite puis éprouvée ; la dernière colonne en délimite la portée.}
\label{tab:apports-institut}
\end{table}
